% Part: incompleteness
% Chapter: theories-computability
% Section: complete-decidable

\documentclass[../../../include/open-logic-section]{subfiles}

\begin{document}

\olfileid{inc}{tcp}{cdc}

\olsection{النظريات التامة \printtoken{S}{axiomatizable} قابلة للقرار}

تسمى النظرية {\em تامة} إذا ثبت فيها، لكل جملة~$!A$، أحد الأمرين:
$!A$ أو $\lnot !A$.

\begin{lem}
كل نظرية $\Th{T}$ تامة و!!{axiomatizable} تكون
$\Th{T}$ قابلة للقرار.
\end{lem}

\begin{proof}
لتكن $\Th{T}$ تامة، ولتكن ${\OLId{latin-looped}{A}}$ مجموعة بديهيات قابلة للحساب لها.
فأما إن كانت $\Th{T}$ غير متسقة، فقابليتها للحساب ظاهرة:
خوارزميتها أن تجيب «نعم» في كل مرة. فيبقى أن نفرض $\Th{T}$
متسقة أيضًا.

وللقرار هل الجملة $!A$ في $\Th{T}$، نجري بحثين متزامنين:
عن !!a{derivation} لـ$!A$ من $\Th{T}$، وعن !!a{derivation}
لـ$\lnot
!A$ من $\Th{T}$. فتمام $\Th{T}$ يضمن العثور على أحدهما،
واتساق $\Th{T}$ يضمن أنه متى وجد !!a{derivation} لـ$\lnot !A$
امتنع وجود !!{derivation} لـ$!A$.

وللبرهان عبارة أخرى. قد ثبت أن $\Th{T}$ !!{c.e.}؛ فتكفي، بالمبرهنة
المتقدمة، قابلية متممة $\Th{T}$ للتعداد كذلك، أي كونها
!!{c.e.}. ذلك أن !!a{sentence} $!A$ تنتمي إلى $\Th{\bar T}$ إذا
وفقط إذا انتمت $\lnot !A$ إلى $\Th{T}$؛ وعليه
$\Th{\bar T}\leq_{\OLId{latin-ordinary}{m}}\Th{T}$.
\end{proof}

\end{document}
